% !TEX root = ../main.tex
% =============================================================================
%  5. Results
% =============================================================================
% PROVISIONAL NUMBERS (2026-11-16). Every number below comes from the global solution on the
% canonical calibration of 2026-11-13: Bohn coefficient 0.15, headline shock 2% a quarter,
% 1.5% kept as the lower-stress case. The only policy experiment is the announced proportional
% rule at k = 2 (branch new-rules); the surprise 200 bp cap is archived in
% Claude files/docs/shelved/05-results_with_surprise_cap_2026-11-16.tex.
% Responses clear the impact quarter exactly with next quarter's values from the fitted decision
% rules; numbers in parentheses or quoted as "two-quarter" clear the impact and following quarter
% jointly. Sources: results/GLOBAL_2026-11-13_bohn015 (headline.json, no_policy/run.log,
% forensic/forensic_out.json, forensic/forensic3_out.json).
\section{Results}
\label{sec:results}

This section answers three questions. How does a rise in the priced default probability
$\pi^d_t$ reach output? What does an announced purchase rule do to that transmission? Who bears
its cost? All responses start from the stochastic rest point and are measured against the path
without the shock.

\subsection{State-dependent transmission without a backstop}
\label{sec:res-rest}

Sovereign risk reaches output through bank balance sheets, and only once losses on the bond book
bind the banks' incentive constraint. Below that point banks absorb the losses out of net worth;
above it, the same losses raise the cost of credit. This subsection locates the threshold and
traces the transmission at the headline shock.

At the rest point $D$ banks are better capitalised than at the deterministic steady state. Their
net worth $n_{D,t}$ is 8.0\% higher, because the bonds they hold earn a risk premium that accrues
as retained earnings. The constraint is slack: its multiplier $\mu_{D,t}$ (the shadow cost of the
constraint) is zero, and so is the lending spread $r^{wc}_{D,t}-r_{D,t}$, the premium firms pay
over the deposit rate on working capital.

The constraint binds once the quarterly default probability passes 1.3\%.
\Cref{tab:state-dependence} shows how transmission changes at that point. Below it, a riskier
sovereign lowers the bond price and output, but the lending spread does not move. Above it, each
further point of default probability raises the multiplier and the lending spread, and output
falls faster. Moving $\pi^d_t$ from 1.0\% to 2.0\% raises the fall in output three and a half
times.

\begin{table}[htbp]
\centering
\caption{Impact responses by priced default probability, no backstop}
\label{tab:state-dependence}
\small
\setlength{\tabcolsep}{4pt}
\begin{tabular}{lcccccc}
\toprule
Quarterly default probability $\pi^d_t$ (\%) & 1.0 & 1.5 & 2.0 & 3.0 & 4.0 & 6.0 \\
\midrule
$D$ bond price $q^D_t$ (\%) & $-4.5$ & $-6.2$ & $-7.9$ & $-10.7$ & $-13.2$ & $-17.5$ \\
Multiplier $\mu_{D,t}$ & 0 & 0.008 & 0.025 & 0.052 & 0.072 & 0.102 \\
Lending spread $r^{wc}_{D,t}-r_{D,t}$ (bp) & 0 & 35 & 115 & 233 & 319 & 438 \\
$D$ output (\%) & $-0.029$ & $-0.055$ & $-0.103$ & $-0.174$ & $-0.224$ & $-0.292$ \\
\bottomrule
\end{tabular}
\par\medskip
\begin{minipage}{0.9\linewidth}\footnotesize
\emph{Notes:} impact of a one-off rise in $\pi^d_t$ from its rest value of 0.10\%, all other states
at the stochastic rest point. Prices and output in percent deviation from the rest point; the
lending spread annualised; the multiplier in levels (zero at rest).
\end{minipage}
\end{table}

The threshold lies far in the tail of the risk process. The risk factor $s_t$ has persistence
$\rho_s=0.9$, innovation volatility $\sigma_s=0.41$ and an unconditional standard deviation of
0.94, and the 95th percentile of $\pi^d_t$ is 0.47\% a quarter. The headline shock of 2\% a
quarter (7.8\% a year) lies 3.2 unconditional standard deviations above the mean of $s_t$ (7.4
innovations), and the stationary distribution puts 0.07\% of quarters above it. The lower-stress
shock of 1.5\% lies 2.9 standard deviations above the mean (0.2\% of quarters). The
balance-sheet channel is a crisis phenomenon.

\paragraph{Transmission at the headline shock.}
\label{sec:res-transmission}
The $D$ bond price falls 7.9\% and the sovereign spread rises from 62 to 269\,bp. Banks hold the
bond at market value, so their net worth falls 12.5\%. That is enough to bind the constraint:
$\mu_{D,t}$ rises from zero to 0.025. By \eqref{eq:foc-general}, every asset must then earn an
excess return of $\lambda_D\mu_{D,t}$ over deposits, where $\lambda_D$ is the share of assets a
banker could divert. We call $\lambda_D\mu_{D,t}$ the constraint premium. Firms pay it on working
capital, and the lending spread rises 115\,bp. Investment falls 1.8\% and output 0.103\% on
impact (0.41\% annualised).

The higher lending spread lowers output by 0.087 percentage points through the cost of working
capital. Weaker demand for the $D$ good lowers it by a further 0.060 points through the terms of
trade. The $D$ deposit rate falls 57\,bp, which offsets half of the rise in the lending spread
and adds back 0.043 points.

The bond's expected excess return rises 381\,bp. Following \citet{bocola2016pass}, 267\,bp is a
risk premium: the bond pays least when the banker values wealth most. The other 115\,bp is the
constraint premium, which \citet{bocola2016pass} calls the liquidity premium, so the constraint
accounts for about 30\% of the rise.

The constraint binds only on impact. From the next quarter the lending spread is zero, yet output
stays about 0.047\% below its no-shock path for two years and 0.035\% below after six: the
investment collapse has run down a capital stock that recovers slowly. Net worth overshoots.
Banks bought the bond cheap, and its recovery as risk recedes leaves their net worth 3.3\% above
the no-shock path after two years.

\subsection{The announced proportional backstop}
\label{sec:res-announced}

An announced rule can matter twice: it can shift the economy before any crisis, and it changes
how the economy responds once a crisis hits. This subsection measures both. The first effect is
small; the second is large for credit and small for the bond market.

The central bank follows a purchase rule that every agent knows. It buys $D$ bonds whenever their
expected excess return exceeds its value at the rest point,
\begin{equation}
q^{D}_{t}\,M_{t+1}=k\,\bar q^{D}\bar B_{D}\,\bigl(\mathit{ER}_{t}-\overline{\mathit{ER}}\bigr),
\qquad M_{t+1}\ge(1-\delta^b_D)\,h_t\,M_{t},
\label{eq:proportional-rule}
\end{equation}
where $M_{t+1}$ is its holding of $D$ bonds, $\mathit{ER}_t$ the quarterly expected excess
return of the $D$ bond over the $D$ deposit rate, $\overline{\mathit{ER}}=15$\,bp a year its
value at the rest point without a backstop, $\bar q^{D}\bar B_D$ the market value of the $D$ debt
at the steady state, and $k$ the response coefficient. The central bank never sells and holds its
bonds to maturity. Because the rule responds to the excess return rather than to the spread, it
does not lean against the part of the spread that compensates for expected default losses. We
set $k=2$.

The announcement effect compares the rest points of the economies with and without the rule. The
incremental crisis effect compares their responses to the same shock, each from its own rest
point. Agents in both economies know their policy throughout, so neither comparison is a
surprise. The two effects sum to the total difference between the shocked economies.

The announcement effect is small. At the rest point the central bank holds 0.02\% of the debt,
the sovereign spread is 1.7\,bp lower, investment is 0.10\% higher and output 0.04\% higher. Net
worth is unchanged to within 0.02\%, and the constraint is slack in both economies.

The crisis effect is large. At the headline shock the rule buys 1.4\% of the $D$ debt and lowers
the sovereign spread by only 13\,bp, yet it removes 110--115\,bp of the 115\,bp rise in the
lending spread (\cref{tab:headline-impact}). $D$ bank net worth falls 9.6\% instead of 12.5\%,
the fall in investment halves, and the impact fall in output shrinks from 0.103\% to 0.024\%. At
the 1.5\% shock the rule buys 1.1\% of the debt and removes the whole 35\,bp rise in the lending
spread.

\begin{table}[htbp]
\centering
\caption{Impact of a sovereign-risk shock: no backstop and the announced proportional rule}
\label{tab:headline-impact}
\small
\begin{tabular}{lccc}
\toprule
 & No backstop & Announced rule & Effect of the rule \\
\midrule
\multicolumn{4}{l}{\emph{Default probability 2\% a quarter}} \\
Sovereign spread (bp) & +207 (+202) & +194 (+192) & $-$13 ($-$11) \\
$D$ bond price $q^D_t$ (\%) & $-$7.87 ($-$7.70) & $-$7.42 ($-$7.33) & +0.45 (+0.37) \\
$D$ bank net worth $n_{D,t}$ (\%) & $-$12.5 ($-$11.4) & $-$9.6 ($-$8.9) & +2.9 (+2.5) \\
Multiplier $\mu_{D,t}$ & +0.025 (+0.025) & 0.000 (+0.002) & $-$0.025 ($-$0.024) \\
Lending spread $r^{wc}_{D,t}-r_{D,t}$ (bp) & +115 (+119) & 0 (+8) & $-$115 ($-$110) \\
Investment $I_{D,t}$ (\%) & $-$1.77 ($-$1.45) & $-$0.93 ($-$0.73) & +0.84 (+0.73) \\
Output $Y_{D,t}$ (\%) & $-$0.103 ($-$0.100) & $-$0.024 ($-$0.027) & +0.079 (+0.073) \\
Purchases (\% of s.s.\ $D$ debt) & 0.00 (0.00) & +1.42 (+1.49) & +1.42 (+1.49) \\
Consumption $C_{D,t}/C_{F,t}$ (\%) & +0.096 (+0.063) & +0.105 (+0.080) & +0.009 (+0.017) \\
\midrule
\multicolumn{4}{l}{\emph{Default probability 1.5\% a quarter}} \\
Sovereign spread (bp) & +161 (+157) & +154 (+151) & $-$8 ($-$6) \\
$D$ bond price $q^D_t$ (\%) & $-$6.23 ($-$6.08) & $-$5.96 ($-$5.86) & +0.27 (+0.22) \\
$D$ bank net worth $n_{D,t}$ (\%) & $-$9.8 ($-$8.8) & $-$7.8 ($-$7.2) & +2.0 (+1.6) \\
Multiplier $\mu_{D,t}$ & +0.008 (+0.008) & 0.000 (0.000) & $-$0.008 ($-$0.008) \\
Lending spread $r^{wc}_{D,t}-r_{D,t}$ (bp) & +35 (+37) & 0 (0) & $-$35 ($-$37) \\
Investment $I_{D,t}$ (\%) & $-$1.36 ($-$1.08) & $-$0.78 ($-$0.60) & +0.58 (+0.48) \\
Output $Y_{D,t}$ (\%) & $-$0.055 ($-$0.052) & $-$0.020 ($-$0.020) & +0.035 (+0.033) \\
Purchases (\% of s.s.\ $D$ debt) & 0.00 (0.00) & +1.06 (+1.06) & +1.06 (+1.06) \\
Consumption $C_{D,t}/C_{F,t}$ (\%) & +0.123 (+0.094) & +0.092 (+0.074) & $-$0.031 ($-$0.020) \\
\midrule
\multicolumn{4}{l}{\emph{Announcement effect at the rest point (rule against no backstop)}} \\
Sovereign spread (bp) & & & $-$1.7 ($-$1.7) \\
Output $Y_{D,t}$ (\%) & & & +0.038 (+0.038) \\
Investment $I_{D,t}$ (\%) & & & +0.10 (+0.10) \\
$D$ bank net worth $n_{D,t}$ (\%) & & & $-$0.02 ($-$0.01) \\
\bottomrule
\end{tabular}
\par\medskip
\begin{minipage}{0.95\linewidth}\footnotesize
\emph{Notes:} impact quarter, each economy measured from its own stochastic rest point;
quantities and prices in percent, rates in bp a year, $\mu_{D,t}$ in levels. The effect of the
rule is the rule's response less the no-backstop response; it excludes the announcement effect,
reported separately. Purchases are the face value bought in the impact quarter over the
steady-state $D$ debt. Each cell clears the impact quarter exactly with next quarter's values
from the fitted decision rules; in parentheses, the impact and the following quarter are cleared
jointly. The constraint binds from a default probability of 1.31\% (1.30\%) without a backstop
and 2.39\% (1.92\%) under the rule.
\end{minipage}
\end{table}

\begin{figure}[htbp]
\centering
\includegraphics[width=\textwidth]{figures/fig_irf_announced_2pct}
\caption{Responses to a rise in the quarterly default probability to 2\%, without a backstop
and under the announced proportional rule. Without a backstop the $D$ banks' constraint binds on
impact, and lending and investment fall; under the rule the constraint is slack or nearly so,
while the sovereign spread barely changes. Lines clear each quarter exactly; hollow markers show
the impact values when two quarters are cleared jointly. Output, investment and net worth in
percent deviation from the no-shock path; spreads in bp a year, the sovereign spread in levels;
purchases in percent of the steady-state $D$ debt.}
\label{fig:irf-announced}
\end{figure}

The effect outlasts the purchase. \Cref{fig:irf-announced} shows that the central bank buys only
in the impact quarter; its holdings then run off, to 0.4\% of the debt after six years. Output
under the rule stays about 0.027\% below its no-shock path for two years, against 0.047\%
without a backstop, and the cumulative output loss over six years falls from 1.12 to 0.64
percent-quarters. Net worth is 2.5\% above the no-shock path after two years, against 3.3\%
without a backstop.

\subsection{Expected capital prices, bank net worth and credit}
\label{sec:res-mechanism}

The rule lowers the lending spread through the price of capital, not through the bonds it buys.
Rearranging \eqref{eq:capital-euler} with \eqref{eq:rk}, the banker prices capital as
\begin{equation}
Q_{D,t}=\frac{\E_t\bigl[\Omega_{D,t,t+1}\bigl(mpk_{D,t+1}+(1-\delta_D)\,Q_{D,t+1}\bigr)\bigr]}
             {\E_t\bigl[\Omega_{D,t,t+1}\bigr]R_{D,t}+\lambda_D\,\mu_{D,t}},
\label{eq:capital-price}
\end{equation}
where $Q_{D,t}$ is the price of installed capital, $mpk_{D,t+1}$ next quarter's marginal product
of capital, $\delta_D$ depreciation, $\Omega_{D,t,t+1}$ the banker's discount factor and $R_{D,t}$
the gross deposit rate. The numerator is the expected resale value of capital plus its dividend;
the denominator contains the constraint premium that also sets the lending spread in
\eqref{eq:rwc}.

\Cref{tab:mechanism} separates the anticipation of purchases from the purchase itself. All four
columns start from a common state, the rest point without a backstop with the default
probability raised to 2\%. The first and last are the two equilibria. The middle two are
controlled counterfactuals that clear the impact quarter under assumptions no announced policy
would produce. In the delayed-purchase counterfactual, agents know the rule, but the central bank
buys nothing in the impact quarter and applies the rule from the next quarter on. In the
surprise-purchase counterfactual, the central bank makes the rule's purchase, 1.43\% of the
debt, while agents keep the expectations of the economy without a backstop.

\begin{table}[htbp]
\centering
\caption{Anticipated and current purchases at the headline shock: equilibria and controlled
counterfactuals}
\label{tab:mechanism}
\footnotesize
\setlength{\tabcolsep}{2.5pt}
\begin{tabular}{lcccc}
\toprule
 & No backstop & Surprise purchase & Rule, purchases & Announced rule \\
 & (equilibrium) & (counterfactual) & from next quarter & (equilibrium) \\
 & & & (counterfactual) & \\
\midrule
Purchase (\% of s.s.\ $D$ debt) & 0 & 1.43 & 0 & 1.43 \\
Lending spread (bp) & 115 (119) & 118 (114) & 0 (0) & 0 (0) \\
Price of capital $Q_{D,t}$ (\%) & & +0.015 (+0.02) & +0.443 (+0.40) & +0.458 (+0.40) \\
$D$ bank net worth $n_{D,t}$ (\%) & & +0.15 (+0.17) & +3.4 (+3.1) & +3.6 (+3.1) \\
\bottomrule
\end{tabular}
\par\medskip
\begin{minipage}{0.95\linewidth}\footnotesize
\emph{Notes:} impact quarter, every column at the same state: the stochastic rest point of the
economy without a backstop, with the default probability raised to 2\% a quarter. The price of
capital and net worth in percent deviation from the no-backstop column. ``Surprise purchase'':
the rule's purchase made unexpectedly, under the decision rules of the economy without a
backstop. ``Rule, purchases from next quarter'': the rule is known but buys nothing in the impact
quarter. Lending spread annualised. In parentheses, the impact and the following quarter are
cleared jointly. \Cref{app:mechanism} gives the full decompositions.
\end{minipage}
\end{table}

Anticipated purchases do almost all the work. In the delayed-purchase counterfactual the price of
capital rises 0.44\%, 97\% of its rise under the rule, and the lending spread falls to zero. The
surprise purchase raises the price of capital by 0.015\% and moves the lending spread by a few
basis points, upward with one quarter cleared and downward with two, so its own effect cannot be
told from zero. The purchase is small relative to what it would have to move: 1.4\% of a bond
book that is itself less than a fifth of the assets banks could divert.

Equation \eqref{eq:capital-price} distinguishes the source of the price rise from its
amplification. Of the 0.46\% rise under the rule, 0.28 points come from a higher expected resale
value of capital: banks expect the rule to support the economy in high-risk states, hence more
investment and stronger demand next quarter, which raises next quarter's price of capital
through the adjustment cost. A further 0.29 points come from the fall of the constraint premium
to zero. That term is an equilibrium outcome, not a separate cause: it equals the no-backstop
lending spread per quarter (115\,bp divided by four), and it disappears because the constraint
goes slack. The higher deposit rate needed to finance the extra investment subtracts 0.11 points,
and dividends contribute nothing. These shares are an exact accounting identity at each
allocation.

The constraint goes slack because the price rise revalues capital the banks already hold. Banks
carry the undepreciated stock into the quarter and value it at $Q_{D,t}$, so
\begin{equation}
\Delta n_{D,t}=(1-f_D)(1-\delta_D)\,K_{D,t}\,\Delta Q_{D,t},
\label{eq:revaluation}
\end{equation}
where $f_D$ is the share of bankers who exit each quarter and $K_{D,t}$ the capital stock.
Capital is 6.6 times net worth at the headline shock, so the 0.46\% price rise raises net worth
by 3.0\% of its no-backstop value; the capital gain on the sovereign book adds 0.6\%. At the
no-backstop allocation a 1\% rise in the price of capital lowers $\mu_{D,t}$ by about 0.057 (the
semi-elasticity $\partial\mu_{D,t}/\partial\ln Q_{D,t}\approx-5.7$), because net worth rises 6.6
times as fast as the price while the divertible base, mostly capital financed with deposits,
rises only 0.8 times as fast. The 0.46\% rise is worth about 0.026, slightly more than the
no-backstop multiplier of 0.025, so the constraint goes slack. The loop between price and net
worth amplifies the initial impulse: holding the impact-quarter allocation at its no-backstop
value, the change in expectations alone raises the price banks will pay for capital by 0.34\%, a
controlled counterfactual; in equilibrium the rise is 1.36 times as large.

The common-state effect on net worth, 3.6\%, is larger than the 2.9 points of
\cref{tab:headline-impact} for two reasons. The policy economy starts from its own rest point,
where banks carry more capital financed with deposits and slightly fewer bonds, which lowers the
effect to 3.3\% of shocked no-backstop net worth. And \cref{tab:headline-impact} expresses
responses relative to rest-point net worth, which is 14\% larger than shocked net worth. Clearing
two quarters jointly lowers the price and net-worth effects by about an eighth and leaves the
ranking of sources unchanged. At the 1.5\% shock the price of capital rises 0.32\% under the
rule, 0.27 points of it from the expected resale value, and the delayed-purchase counterfactual
delivers 96\% of the rise.


\subsection{Cross-country incidence and limitations}
\label{sec:res-incidence}

Two questions remain: who bears the cost of the rule, and how far the numbers can be trusted.
The rule moves little risk onto $F$. Its effect on credit survives a better approximation; the
exact level of the lending spread it leaves behind does not.

Without a backstop $F$ is almost untouched. $F$ banks hold 12.7\% of the $D$ debt, but $F$ is
11.3 times larger than $D$, so the loss per unit of $F$ bank capital is negligible. At the
headline shock $F$ output and bank net worth even rise slightly on impact, by 0.007\% and
0.07\%. The rule changes little. Its holdings earn the central bank an expected profit next
quarter of 0.02\% of $D$'s quarterly output, of which $D$ receives the capital-key share
$\kappa=0.071$; if default occurs, the loss is 1.2\% of $D$'s quarterly output. $D$ consumption
relative to $F$ consumption moves by at most 0.03 percentage points on impact. A rule that works
through anticipation places little risk on the common balance sheet.

The accounting holds exactly at every state we report, under both clearing methods. Bank net
worth equals the returns on its assets less its deposits, the $D$ bond market clears, $\mu_{D,t}$
is the positive part of the constraint gap, the lending spread equals the shadow cost of the
constraint, and \eqref{eq:capital-price} holds, all to within $10^{-10}$. The central bank's
holdings are never negative along any reported response or in a 1,200-quarter simulation.

The size of the policy effect depends on the approximation. Each response clears the impact
quarter exactly but takes next quarter's values from the fitted decision rules, which overstate
next quarter's value of bank net worth by up to 3.4\% at the headline shock. Clearing two
quarters jointly changes the rule's effect on the lending spread from 115 to 110\,bp and on
output from 0.079 to 0.073 points. We take that 5\% gap as the measure of approximation error
near the constraint and claim no convergence beyond it.

The lending spread the rule leaves at 2\% is therefore not established. It is zero with one
quarter cleared. With two quarters cleared it is 8\,bp from the policy economy's own rest point
($\mu_{D,t}=0.002$) but zero from the common state of \cref{tab:mechanism}. The headline shock
sits where the rule's constraint starts to bind, at 2.39\% with one quarter cleared and 1.92\%
with two, so whether the constraint is exactly slack or barely binding depends on the
approximation and on the initial state. Under either method the rule removes at least 110 of the
115\,bp. At higher risk it no longer keeps the constraint slack, leaving lending spreads of
7\,bp at 2.5\% and 37\,bp at 3\%.
